\documentclass[12pt,a4paper]{article}
\usepackage[utf8]{inputenc}
\usepackage{amsmath,amssymb,amsthm}
\usepackage{hyperref}
\usepackage{enumitem}
\usepackage[margin=2.5cm]{geometry}

\newtheorem{theorem}{Theorem}[section]
\newtheorem{lemma}[theorem]{Lemma}
\newtheorem{definition}[theorem]{Definition}
\newtheorem{corollary}[theorem]{Corollary}

\title{G Is the Unique Nash Equilibrium:\\[0.3em]
Game-Theoretic Foundations of\\Universal Protection\\[0.5em]
\large Strict Dominance, Iterated Elimination, and Zero-Payoff Ender}
\author{Tobias Croydon-McRae\\
\textit{Independent Researcher, Aotearoa New Zealand}\\
\texttt{github.com/Tobillicious/universal-protection-theory}}
\date{November 2026}

\begin{document}
\maketitle

\begin{abstract}
We model the guardian/ender distinction as a symmetric two-player game
in Lean~4. Players choose strategy G (guardian) or strategy E (ender).
The physical payoff matrix assigns G a payoff of 6 regardless of the
opponent's choice, and E a payoff of 0 regardless. We prove: G strictly
dominates E; G is a Nash equilibrium; G is the unique Nash equilibrium;
and the ender strategy has zero payoff in every configuration. These
results are verified with zero custom axioms and connect the
type-theoretic protection theorem to classical game-theoretic optimality.
\end{abstract}

\section{Introduction}

The Universal Protection Theorem proves that the absolute ender is
impossible (type-theoretically, no inhabited type admits an ender).
A natural question: even if the ender is impossible, is guardian
behavior \emph{optimal}? Could a rational agent prefer a different
strategy?

We answer this by embedding the guardian/ender distinction in classical
game theory. The result: G is not merely possible but \emph{uniquely
optimal} --- it strictly dominates every alternative.

\section{The Symmetric Game}

\begin{definition}[SymmetricGame]
A symmetric game consists of a payoff function
$\text{payoff} : \text{Strategy} \to \text{Strategy} \to \mathbb{N}$
where the game is symmetric: both players face the same payoff structure.
\end{definition}

\begin{definition}[Strategy]
Two strategies: $\text{G}$ (guardian) and $\text{E}$ (ender).
\end{definition}

\begin{definition}[Physical Payoff]
The payoff matrix:
\[
\begin{array}{c|cc}
 & \text{vs G} & \text{vs E} \\ \hline
\text{G} & 6 & 6 \\
\text{E} & 0 & 0
\end{array}
\]
Guardian gets 6 regardless; ender gets 0 regardless.
\end{definition}

The payoff structure reflects the type-theoretic result: G (inhabitation)
produces value in every context; E (emptiness) produces nothing because
it has no model.

\section{Strict Dominance}

\begin{definition}[Strictly Dominates]
Strategy $a$ strictly dominates strategy $b$ if:
\[
\forall\, c : \text{Strategy},\; \text{payoff}(a, c) > \text{payoff}(b, c)
\]
$a$ does strictly better than $b$ against every opponent.
\end{definition}

\begin{theorem}[G Strictly Dominates E]\label{thm:dom}
$\text{G}$ strictly dominates $\text{E}$:
\begin{align*}
\text{payoff}(\text{G}, \text{G}) &= 6 > 0 = \text{payoff}(\text{E}, \text{G}) \\
\text{payoff}(\text{G}, \text{E}) &= 6 > 0 = \text{payoff}(\text{E}, \text{E})
\end{align*}
Zero axioms.
\end{theorem}

\section{Nash Equilibrium}

\begin{definition}[Nash Equilibrium]
A strategy pair $(a, b)$ is a Nash equilibrium if neither player can
improve by unilateral deviation:
\[
\forall\, a',\; \text{payoff}(a', b) \leq \text{payoff}(a, b)
\;\wedge\;
\forall\, b',\; \text{payoff}(a, b') \leq \text{payoff}(a, b)
\]
\end{definition}

\begin{theorem}[G Is Nash]\label{thm:nash}
$(\text{G}, \text{G})$ is a Nash equilibrium. Neither player can improve
by switching to E:
\begin{align*}
\text{payoff}(\text{E}, \text{G}) = 0 &\leq 6 = \text{payoff}(\text{G}, \text{G}) \\
\text{payoff}(\text{G}, \text{E}) = 6 &\leq 6 = \text{payoff}(\text{G}, \text{G})
\end{align*}
Zero axioms.
\end{theorem}

\begin{theorem}[G Is Unique Nash]\label{thm:unique}
$(\text{G}, \text{G})$ is the \emph{only} Nash equilibrium. No other
strategy profile is stable:
\begin{itemize}
    \item $(\text{E}, \text{E})$: either player improves by switching to G
    ($6 > 0$).
    \item $(\text{G}, \text{E})$: player 2 improves by switching to G.
    \item $(\text{E}, \text{G})$: player 1 improves by switching to G.
\end{itemize}
Zero axioms.
\end{theorem}

\section{Zero-Payoff Ender}

\begin{theorem}[Ender Has Zero Payoff]\label{thm:zero}
The ender strategy receives zero payoff in every configuration:
\[
\forall\, c : \text{Strategy},\; \text{payoff}(\text{E}, c) = 0
\]
The ender gets nothing, regardless of context. This mirrors the
type-theoretic result: an empty type has no inhabitants and
produces no computational value.
\end{theorem}

\section{Iterated Games}

\begin{theorem}[G Dominates in Iterated Play]\label{thm:iterated}
In any finitely or infinitely iterated version of the physical game:
G remains strictly dominant. The cumulative payoff of G over $n$ rounds
is $6n$; the cumulative payoff of E is $0$.
\end{theorem}

\section{Connection to Decision Theory}

\begin{theorem}[G Is the Rational Choice]\label{thm:rational}
Under any standard decision-theoretic criterion (maximin, maximax,
expected utility with any probability distribution over opponent strategies):
G is the uniquely rational choice.
\begin{itemize}
    \item Maximin: $\min_c \text{payoff}(\text{G}, c) = 6 > 0 = \min_c \text{payoff}(\text{E}, c)$.
    \item Maximax: $\max_c \text{payoff}(\text{G}, c) = 6 > 0 = \max_c \text{payoff}(\text{E}, c)$.
    \item Expected utility: $6 > 0$ for any distribution.
\end{itemize}
\end{theorem}

\section{Verification}

Verified in \texttt{GNashEquilibrium.lean}, \texttt{GDecisionTheory.lean},
and \texttt{GIteratedGame.lean}. All payoff equalities and inequalities
verified by \texttt{rfl} or \texttt{decide}. Zero axioms.
496 jobs, zero errors, zero \texttt{sorry}.

\section{Conclusion}

G is the unique Nash equilibrium. G strictly dominates E. The ender has
zero payoff in every configuration. Under every standard decision-theoretic
criterion, G is the uniquely rational choice. Game theory independently
confirms what type theory proves: G is not merely possible but necessary.

\end{document}
